% numerical_cross_validation.tex -- 可复核数值证据与独立交叉验证
\section{数值交叉验证与可复现实验}\label{sec:gr-numerical-crossval}

\subsection{实验边界与复现条件}

本节补上实现审计不能由断言数量替代的数值证据。所有数值来自同一轮当前工作树
Release 重编译后的注册测试。实验机为 Apple Silicon
macOS/arm64，CMake preset 为 \srcpath{macos-release}（Release、\texttt{-O3}、
\texttt{-DNDEBUG}）。主仓库基准 HEAD 为 \texttt{114f3225f4c3} 加本轮未提交变更，
本地 \texttt{../MIPSolvers} 为且等于仓库 pin
\texttt{4a0b16a00daefcbe18d2c5648fe30f840ed77052}，依赖工作树 clean。

重建和运行命令如下：
\begin{verbatim}
cmake --build --preset macos-release --target \
  test_graph test_graph_kron test_graph_roundtrip \
  test_topology_crossval test_validation run_gui_server -j4
build/macos-release/tests/test_graph
build/macos-release/tests/test_graph_kron
build/macos-release/tests/test_graph_roundtrip
build/macos-release/tests/test_topology_crossval
build/macos-release/tests/test_validation
ctest --test-dir build/macos-release -R gui_api_e2e
\end{verbatim}

focused Release 结果为：\texttt{test\_graph} 37 cases/193 assertions，
\texttt{test\_graph\_kron} 10/56，\texttt{test\_graph\_roundtrip} 17/294，
\texttt{test\_topology\_crossval} 12/86，\texttt{test\_validation} 38/180；合计
114/809。前四个 graph 二进制串行约 0.51 s，低于 5 s 门槛。GUI E2E 的 80 项检查全部通过。
定向 ASan/UBSan CTest 通过 42 项，1 项按既有条件 skip，无 sanitizer 报告。没有运行完整 CTest 或
图模块级全网络规模基准，故不得将本节表格解释为全套认证。

\subsection{验收门槛}

\begin{center}\small
\begin{tabular}{@{}L{5.0cm}L{3.0cm}L{5.2cm}@{}}
\toprule
比较对象 & 验收门槛 & 口径\\
\midrule
无源线性 Kron & boundary current 与恢复状态最大绝对误差 $\le 10^{-10}$ & 复数电流/电压，标幺或导纳单位按测试输入\\
闭合开关、无源 series & 保留电压与总有功损耗差 $\le 10^{-6}$ & 与完整网络 PF 的同一 MW 基准比较\\
case33bw split-series & 恢复电压最大误差 $\le 10^{-4}\,{\rm pu}$，损耗差 $\le 10^{-4}\,{\rm MW}$ & 允许 MATPOWER/数值线性代数舍入\\
DC OPF round-trip & 目标差 $\le 10^{-3}$ & 目标值不是物理损耗\\
\bottomrule
\end{tabular}
\end{center}

\subsection{稠密 Kron：解析式与完整矩阵对照}

三母线无源测试将中间母线消去。稠密 Schur 补得到等值导纳虚部
\(6.66666666666666785\)，解析值为 \(6.66666666666666696\)，误差小于
\(10^{-10}\)。加入边界电压后，边界电压检查为 0（容差 \(10^{-10}\)）。这是一条
独立于实现矩阵填充顺序的标量解析对照，证明无源内部节点的 Kron 路径没有把
导纳符号或端点顺序写反。

对稀疏 Kron tape，保留 3 个节点、消去 2 个节点，并同时用完整矩阵直接求解和
recovery operator 回放。三组边界电流最大误差分别为
\(3.55\times10^{-15}\)、\(2.03\times10^{-15}\)、\(6.30\times10^{-16}\)；
消去节点零注入残差为 \(8.90\times10^{-16}\) 与 \(4.19\times10^{-15}\)，
恢复算子最大误差为 0。该结果满足比 \(10^{-10}\) 更严格的门槛，且同时检查了
边界方程、内部零注入和状态恢复三种量。

新增常电流注入回归同时检查 reduced boundary equation 与完整分块回代，内部节点残差
不超过 \(10^{-12}\)。这证明保存的 $Y_{\beta\beta}^{-1}I_\beta$ correction 对线性常电流模型
闭环；恒功率负荷若只在工作点转换为电流仍不能被描述为全局精确 Ward 等值，调用方必须保留
工作点、注入假设和后验 PF 结果。

\subsection{开关收缩与串联化简的 PF round-trip}

闭合 AC 开关收缩后，两个被合并母线的电压分别为 1.0 pu 与
0.9898411045193044 pu；完整网和收缩网回放的总有功损耗均为
0.01479915889694494 MW，差值低于 \(10^{-6}\,\rm MW\)。该数值也暴露了工程
边界：实现把低阻/闭合关系视为可收缩候选，但输入若允许明显电压差，收缩前后的
电压并非“强制相等”的物理测量值，仍需后验检查。

对被动三母线串联支路，完整/化简 PF 损耗和恢复电压为
\[
\begin{aligned}
P_{\rm loss}^{full}&=0.01430375711174059\ {\rm MW},&
P_{\rm loss}^{red}&=0.01430375711175214\ {\rm MW},\\
V_{2}^{rec}&=0.99381398947056943\ {\rm pu},&
V_{2}^{full}&=0.99381398947057409\ {\rm pu}.
\end{aligned}
\]
损耗差约 \(1.16\times10^{-14}\,\rm MW\)。
把 branch component index 改成非连续值后得到相同数量级结果，说明测试覆盖了
“稳定组件 ID 不等于 vector position”的真实接口风险。

IEEE case33bw 的 split-series round-trip 进一步把多个串联记录组合起来：共生成
3 条 series records，最大恢复电压误差约 \(1.17\times10^{-12}\,\rm pu\)，代表值为
\(0.9985161219139157\) 对 \(0.99851612191391848\)、
\(0.94375047987776928\) 对 \(0.94375047987795657\)、
\(0.91963815539739369\) 对 \(0.91963815539855953\)。完整与化简 PF 的损耗
均为 0.2026771169693 MW 量级，差满足 \(10^{-4}\,\rm MW\) 门槛。

\subsection{DC OPF 与图化简的一致性}

串联化简前后的 DC OPF 目标分别为
\(800.0\) 与 \(800.00000000000011369\)，差约
\(1.14\times10^{-13}\)。这只是目标函数 round-trip 证据：它不证明 AC PF
可行性，也不证明经济成本与物理损耗相同；后者仍须由 OPF/PF 各自的 result scope
和 residual 字段给出。

\subsection{证据解读与未覆盖范围}

上述结果形成四层交叉验证：解析三母线标量、完整矩阵与 recovery operator、完整/化简网络的
PF/OPF round-trip，以及富组件 HTTP export/reload。AUD-024--AUD-031 的缺陷路径均已有定向动态
回归并按声明范围关闭。仍未覆盖任意富组件的全网络电气等值基准；Transformer3W/LCC/
EnergyRouter 图边只表示连接性，HTTP Kron identify-only，pendant 为工作点近似。因此生产使用仍须
保存域限定 ID、原系统及后验 PF/OPF 证据。
